\section{Erasmus+ Prácticas}

% 2. ¿Qué es Erasmus+ Prácticas?
\begin{frame}{¿Qué es Erasmus+ Prácticas?}
    \begin{itemize}
        \item Movilidad para realizar prácticas en empresas, institutos de investigación o laboratorios en el extranjero.
        \item \textbf{Objetivo:} Adaptarse al mercado laboral europeo y adquirir experiencia profesional internacional.
      \item \textbf{Grado:} Haber superado al menos el 50\% de los créditos.
      \item \textbf{Máster:} Tener superadas al menos 2 asignaturas.
      \item \textbf{Idiomas:} Poseer conocimiento suficiente de la lengua del país de destino.
      \item \textbf{Remuneración de la empresa:} La ayuda Erasmus+ \textbf{es compatible} con cualquier salario o ayuda económica que la institución de acogida decida abonar al estudiante.        
    \end{itemize}
\end{frame}

% Info relevante
\begin{frame}{Info Relevante}
    \begin{itemize}

      \item \textbf{Duración:} 2 a 12 meses por ciclo de estudios (Grado, Máster o Doctorado).
      \item \textbf{Duración máxima financiable:} 4 meses.
      \item \textbf{Fecha límite:} Todas las prácticas deben finalizar antes del \textbf{31 de julio de 2027}.
        \item El procedimiento de reconocimiento de créditos debe realizarse \textbf{antes} de solicitar la movilidad.
        \item Es obligatorio contar con el visto bueno del tutor académico de la URJC y del tutor de la institución de acogida.
    \end{itemize}
\end{frame}


% 6. Pasos Clave del Estudiante
\begin{frame}{¿Cómo empezar?}
    Responsabilidades del estudiante según la convocatoria:
    \begin{enumerate}
        \item \textbf{Buscar la institución:} Empresa o centro que encaje con tu perfil.
        \item \textbf{Recursos Específicos para Erasmus+:}
            \begin{itemize}
                \item \textbf{Erasmus Intern:} Portal de la red ESN (\url{erasmusintern.org}) específico para prácticas.
                \item \textbf{EURES:} El portal europeo de la movilidad profesional.
            \end{itemize}
        \item \textbf{Ser seleccionado:} La institución de acogida debe aceptarte.
        \item \textbf{Visto bueno:} Obtener la aprobación del tutor académico de la URJC.
        \item \textbf{Acuerdo de Formación:} Firmar el \textit{Training Agreement} (Anexo III) antes de marcharte.
    \end{enumerate}
    \vspace{0.3cm}
    \centering
    \small{Consulta más información en:\\ \url{https://www.urjc.es/erasmus-practicas}}
\end{frame}


% 5. Dotación Económica
\begin{frame}{Ayudas Económicas}
    Las cuantías mensuales dependen del coste de vida del país de destino:
    \begin{table}
        \centering
        \begin{tabular}{|l|c|}
            \hline
            \textbf{Grupo de Países} & \textbf{Importe (Base + Prácticas)} \\ \hline
            Grupo 1 (Vida alta) & 350€ + 150€ = \textbf{500€/mes} \\ \hline
            Grupo 2 (Vida media) & 300€ + 150€ = \textbf{450€/mes} \\ \hline
            Grupo 3 (Vida baja) & 250€ + 150€ = \textbf{400€/mes} \\ \hline
        \end{tabular}
    \end{table}
    \begin{itemize}
        \item \textbf{Nota importante:} La obtención de la plaza no garantiza necesariamente la ayuda financiera (sujeta a fondos).
        \item \textbf{Oportunidades menos favorecidos:} Ayuda adicional de 250€/mes bajo ciertos supuestos (beca Ministerio, familia numerosa...).
    \end{itemize}
\end{frame}

% 8. Clasificación de Países por Grupos
\begin{frame}{Clasificación de Países por Grupos}
    \small
    \begin{table}
        \centering
        \begin{tabular}{|p{1.8cm}|p{8.5cm}|}
            \hline
            \textbf{Grupo} & \textbf{Países de Destino} \\ \hline
            \textbf{Grupo 1} (Vida alta) & Alemania, Austria, Bélgica, Dinamarca, Finlandia, Francia, Irlanda, Islandia, Italia, Liechtenstein, Luxemburgo, Noruega, Países Bajos, Suecia. También: Reino Unido, Suiza, Andorra, Mónaco, San Marino, Vaticano. \\ \hline
            \textbf{Grupo 2} (Vida media) & Chequia, Chipre, Eslovaquia, Eslovenia, España, Estonia, Grecia, Letonia, Malta, Portugal. \\ \hline
            \textbf{Grupo 3} (Vida baja) & Bulgaria, Croacia, Hungría, Lituania, Macedonia del Norte, Polonia, Rumanía, Serbia, Turquía. \\ \hline
        \end{tabular}
    \end{table}
\end{frame}

% 5. Resumen
\begin{frame}{Secuencia Recomendada}

    \begin{enumerate}
        \item Comprobar si eres elegible
        \item Buscar empresa (tutor externo) donde realizar las prácticas
        \item Determinar tareas, duración, condiciones...
        \item Contactar con tutor URJC para visto bueno
        \item Esperar publicación de convocatoria (marzo/abril)
        \item Enviar en {\bf 14 días} documentación necesaria (Anexos II y III)
        \item Esperar a la resolución de la convocatoria
        \item Firmar convenio subvención (Anexo IV)
        \item {\bf Realizar las prácticas Erasmus+} (enviar Anexo V y quizás VI)
        \item Remitir documentación (Anexo VII y documentación prácticas)
    \end{enumerate}
\end{frame}

